\section{Conclusion}
\label{sec5}
In summary, we have introduced a hardware-efficient protocol to dynamically suppress cavity dephasing caused by flux noise inherited from a coupled nonlinear element. By applying a continuous, off-resonant microwave drive to the FTT, we generate a photon-number-dependent AC-Stark shift that cancels the cavity frequency's first-order susceptibility to coupler-frequency fluctuations. Our theoretical analysis and Monte Carlo simulations show that operating at this driven sweet spot can extend the cavity dephasing time by more than an order of magnitude, achieving a ${\sim}20\times$ improvement from ${\sim}100~\mu$s to ${\sim}2$~ms.

While the continuous drive inevitably introduces secondary decoherence channels, such as photon-shot-noise dephasing from drive-induced FTT excitation, we show that operating in the unselective regime at larger detunings parametrically suppresses these effects. In this regime, the primary flux-noise cancellation dominates, resulting in a net enhancement of coherence. Crucially, this mitigation strategy is highly hardware-efficient: it circumvents the need for fast flux tuning during idling, thereby avoiding unwanted transitions that can occur when the coupler frequency is swept past cavity resonances.

Furthermore, we have shown that the mechanism is not limited to the FTT: because it is rooted in drive-induced energy shifts, it extends to other flux-dependent nonlinear elements such as SNAIL-based couplers. We also showed that the protocol is directly relevant for bosonic quantum error correction, where a single drive setting can simultaneously suppress the flux sensitivity of multiple critical transitions in binomial and dual-rail encodings. Overall, the approach provides a practical route toward preserving the noise bias needed for cavity-based QEC at non-sweet-spot flux operating points.